\documentclass[11pt]{article}
\usepackage[letterpaper,margin=22mm]{geometry}
\usepackage{fontspec}
\defaultfontfeatures{Ligatures=TeX}
\setmainfont{Times New Roman}
\setsansfont{Times New Roman}
\setmonofont{Courier New}
\usepackage{amsmath,mathtools,unicode-math}
\setmathfont{STIX Two Math}
\usepackage{graphicx,booktabs,tabularx,xltabular,array,ragged2e}
\graphicspath{{figures/}{supplement_figures/}}
\usepackage{caption,float,placeins,needspace,setspace,microtype}
\usepackage{xcolor}
\usepackage[hidelinks]{hyperref}
\hypersetup{pdftitle={Supplementary Material: Coordinate-Explicit and Route-Specific Finite-Precision Qualification of Tangent-Only Perturbations}}
\renewcommand{\tabularxcolumn}[1]{>{\RaggedRight\arraybackslash}p{#1}}
\renewcommand{\arraystretch}{1.14}
\captionsetup{font=small,labelfont=bf,skip=5pt}
\setlength{\parindent}{1.5em}
\setlength{\parskip}{0.08em}
\setlength{\emergencystretch}{2em}
\setstretch{1.04}
\title{\bfseries Supplementary Material\\[0.4em]\large Coordinate-Explicit and Route-Specific Finite-Precision Qualification of Tangent-Only Perturbations in Coupled Finite-Element Linearizations}
\author{Bing Liang \and Yangqi Ma \and Weiji Sun\thanks{Corresponding author: sunweiji-1231@163.com} \and Shi He}
\date{Computer Methods in Applied Mechanics and Engineering candidate v0.13.0}

\begin{document}
\maketitle
\renewcommand{\thetable}{S\arabic{table}}
\renewcommand{\thefigure}{S\arabic{figure}}

\section*{Purpose, version, and scope}
This supplement accompanies manuscript v0.13.0. It preserves the registered
analytical checks, exact finite-element operand provenance, route-specific
finite-precision error components, response-route validation, independent
precision audit, effectivity anatomy, tooling lineage, and claim boundaries.
The M8 and M16 cases are independent fixed finite-element linearizations, not a
mesh sequence. G2-20 added synthetic response-route evaluation and a newly
generated holdout, but it did not execute Abaqus, modify a finite-element
operator, or import the scalable sparse-estimator results reserved for Paper 3.

\section{Analytical verification}
\begin{table}[htbp]\centering\small
\caption{Registered reproduction of the directional non-determination construction}\label{tab:S1}
\begin{tabularx}{0.98\textwidth}{@{}>{\raggedright\arraybackslash}p{0.07\textwidth}XXXXX@{}}
\toprule
$\varepsilon$ & $\mu_J$ (matched) & $\rho_1$ (matched) & $\rho_g$ (matched) & aligned $\eta_z$ & orthogonal $\eta_z$ \\
\midrule
1e-12 & 1.0000e-12 & 1.0000e-12 & 1.0000e-12 & 1.0000e-12 & 0 \\
1e-10 & 1.0000e-10 & 1.0000e-10 & 1.0000e-10 & 1.0000e-10 & 0 \\
1e-8 & 1.0000e-08 & 1.0000e-08 & 1.0000e-08 & 1.0000e-08 & 0 \\
1e-6 & 1.0000e-06 & 1.0000e-06 & 1.0000e-06 & 1.0000e-06 & 0 \\
1e-4 & 1.0000e-04 & 1.0000e-04 & 9.9990e-05 & 9.9990e-05 & 0 \\
1e-2 & 1.0000e-02 & 1.0000e-02 & 9.9010e-03 & 9.9010e-03 & 0 \\
\bottomrule
\end{tabularx}

\end{table}

\begin{table}[htbp]\centering\small
\caption{Registered coordinate-transformation and redeclaration calculation}\label{tab:S2}
\begin{tabularx}{0.98\textwidth}{@{}>{\raggedright\arraybackslash}p{0.07\textwidth}XXXXX@{}}
\toprule
$s$ & transformed $\widehat{\Delta J}_{22}$ & transformed physical $(\Delta J)_{22}$ & transformed $\eta_z$ & redeclared physical $(\Delta J)_{22}$ & redeclared $\eta_z$ \\
\midrule
1 & 1.0000e-06 & 1.0000e-06 & 1.0000e-06 & 1.0000e-06 & 1.0000e-06 \\
10 & 1.0000e-05 & 1.0000e-06 & 1.0000e-06 & 1.0000e-07 & 1.0000e-07 \\
1e3 & 1.0000e-03 & 1.0000e-06 & 1.0000e-06 & 1.0000e-09 & 1.0000e-09 \\
1e6 & 1.0000e+00 & 1.0000e-06 & 1.0000e-06 & 1.0000e-12 & 1.0000e-12 \\
\bottomrule
\end{tabularx}

\end{table}

\begin{table}[htbp]\centering\small
\caption{Historical reproduction and direct-route coverage for the two reference problems}\label{tab:S3}
\begin{tabularx}{0.98\textwidth}{@{}>{\raggedright\arraybackslash}p{0.24\textwidth}>{\raggedright\arraybackslash}p{0.23\textwidth}>{\raggedright\arraybackslash}p{0.20\textwidth}X@{}}
\toprule
\textbf{Dataset/lane} & \textbf{Coverage} & \textbf{Severe undercoverage} & \textbf{Evidence role} \\
\midrule
historical development & 287/300 (13 uncovered) & 3 & reproduced motivation \\
development source & 300/300 (0 uncovered) & 0 & pre-specified coverage \\
development binary & 300/300 (0 uncovered) & 0 & pre-specified coverage \\
holdout source & 300/300 (0 uncovered) & 0 & pre-specified coverage \\
holdout binary & 300/300 (0 uncovered) & 0 & pre-specified coverage \\
\bottomrule
\end{tabularx}

\end{table}

\begin{figure}[htbp]
\centering
\includegraphics[width=0.94\linewidth,height=0.65\textheight,keepaspectratio]{S1_historical_undercoverage_attribution.pdf}
\caption{Source-representation attribution for the three historical severe-undercoverage cases. This retained diagnostic does not assign a perturbation defect}\label{fig:S1}
\end{figure}

\begin{figure}[htbp]
\centering
\includegraphics[width=0.78\linewidth,height=0.58\textheight,keepaspectratio]{S2_historical_response_label_bins.pdf}
\caption{Historical descriptive response bins retained for provenance only. The labels have no validated class boundary and are excluded from confirmatory conclusions}\label{fig:S2}
\end{figure}

\section{Response-route theorem validation}
For the response equation
\[
  A\delta z=-Ez,\qquad A=J+E,
\]
the computed route uses
\[
  \widetilde f=\operatorname{fl}(-E\widetilde z),\qquad
  r_\delta=\widetilde f-A\widetilde{\delta z},\qquad
  q_f=\widetilde f+E\widetilde z.
\]
With majorants
$\beta_A\geq\lVert A^{-1}\rVert_2$,
$\alpha_E\geq\lVert E\rVert_2$, and
$\lVert\widetilde z-z\rVert_2\leq U_z$, the implemented bound is
\[
 U_{\rm r}
 =
 \beta_A\alpha_EU_z
 +\beta_A\lVert q_f\rVert_2
 +\beta_A\lVert r_\delta\rVert_2.
\]
The development design was reused only after this route model was fixed. A new
300-system holdout then used disjoint seeds and registered condition,
dimension, perturbation-amplitude, and family grids.

\begin{table}[htbp]\centering\footnotesize
\caption{Direct- and response-route coverage and lower-quantity counts}\label{tab:S4}
\begin{tabularx}{\linewidth}{>{\RaggedRight\arraybackslash}p{0.16\linewidth} >{\RaggedRight\arraybackslash}p{0.16\linewidth} *{4}{>{\centering\arraybackslash}X}}
\toprule
Set & Reference contract & Direct coverage & Response coverage & Lower response error & Lower response enclosure \\
\midrule
Development & Decimal source & 300/300 & 300/300 & 295/300 & 283/300 \\
Development & Binary64 operand & 300/300 & 300/300 & 199/300 & 229/300 \\
New holdout & Decimal source & 300/300 & 300/300 & 271/300 & 229/300 \\
New holdout & Binary64 operand & 300/300 & 300/300 & 160/300 & 188/300 \\
\bottomrule
\end{tabularx}

\end{table}

All 2,400 route--contract--split coverage decisions passed: 300/300 for both
routes in both reference contracts on development and on the new holdout.
The response equation had lower decimal-source reference error in 295/300
development and 271/300 new-holdout systems. Binary64 counts were 199/300 and
160/300. These data do not establish a universal route ranking.

\section{Precision and implementation independence}
The registered high-precision implementation was repeated at 120 and 180
decimal digits. The maximum normalized difference was
$5.16\times10^{-110}$ on development and $6.82\times10^{-107}$ on the new
holdout, with no coverage-decision change. A separately implemented
Python \texttt{Decimal} backend evaluated 36 stratified holdout systems in
both reference contracts. Across 72 lane comparisons, the maximum normalized
difference from the registered reference was $6.96\times10^{-153}$.

The first G2-20 execution stopped because its precision comparator normalized
near-zero quantities by their own magnitude. That attempt remains classified
as an inconclusive reference-or-precision result. Two tooling-only corrections
replaced the comparator without changing operators, thresholds, response
formulas, or scientific outputs; the corrected run restarted from row 1. This
lineage is reported rather than overwritten.

\section{Effectivity anatomy and extreme ratios}
For positive reference error $e_{\rm ref}$ and response magnitude
$d_{\rm ref}$, define
\[
 I_{\rm eff}=\frac{U}{e_{\rm ref}},
 \qquad
 R_U=\frac{U}{d_{\rm ref}}.
\]
For a parent bound $U=\lVert A^{-1}\rVert_F\lVert r\rVert_2$ and exact
forward error $e=\lVert A^{-1}r\rVert_2$,
\[
 I_{\rm eff}
 =
 \underbrace{\frac{\lVert A^{-1}\rVert_F}
 {\lVert A^{-1}\rVert_2}}_{\text{norm-majorant inflation}}
 \underbrace{\frac{\lVert A^{-1}\rVert_2\lVert r\rVert_2}
 {\lVert A^{-1}r\rVert_2}}_{\text{directional inflation}}.
\]

\begin{table}[htbp]\centering\footnotesize
\caption{Factorized direct-route effectivity anatomy}\label{tab:S5}
\begin{tabularx}{\linewidth}{>{\RaggedRight\arraybackslash}p{0.18\linewidth} >{\RaggedRight\arraybackslash}p{0.25\linewidth} *{2}{>{\centering\arraybackslash}X}}
\toprule
Set & Quantity & Median range & 95th-percentile range \\
\midrule
Development & Frobenius majorization & $1.0001$--$1.0001$ & $1.94$--$1.94$ \\
Development & Directional inflation & $1.18\times10^{4}$--$2.31\times10^{4}$ &
$3.18\times10^{8}$--$8.53\times10^{8}$ \\
Original holdout & Frobenius majorization & $1.0050$--$1.0050$ & $2.69$--$2.69$ \\
Original holdout & Directional inflation & $1.36\times10^{3}$--$1.88\times10^{3}$ &
$3.59\times10^{7}$--$8.72\times10^{7}$ \\
\bottomrule
\end{tabularx}

\end{table}

The norm-majorant median was approximately $1.0001$ on development and
$1.0050$ on the original holdout. Directional-inflation medians instead
ranged from $1.36\times10^3$ to $2.31\times10^4$. Thus the central
conservativeness is associated primarily with replacing the realized residual
direction by a worst-direction inverse norm.

The largest finite direct-route effectivity was
$7.9321\times10^{86}$, reported in the manuscript as
$\log_{10}I_{\rm eff}=86.90$. It occurred for
\texttt{HOLDOUT\_SPD\_D16\_C01\_E3eM17}, whose reference error was
approximately $4.9995\times10^{-102}$. The extreme ratio is therefore a
near-zero-denominator event, not undercoverage. Exact strings and case
identifiers remain in \texttt{Paper2\_extreme\_ratio\_record.json}.

A registered main-effects analysis gave adjusted $R^2=0.869$ on development;
the drop-one partial $R^2$ for condition exponent was $0.855$. The original
holdout fit gave adjusted $R^2=0.057$ and condition partial $R^2=0.029$.
These finite-grid results support mechanism analysis within the registered
design, not a universal conditioning law.

\section{Frozen finite-element operand provenance}

\begin{tabularx}{\linewidth}{@{}p{0.12\linewidth}p{0.12\linewidth}p{0.24\linewidth}X@{}}
\toprule
\textbf{Case} & \textbf{Dimension} & \textbf{Snapshot SHA-256} & \textbf{Declared operand contract} \\
\midrule
M8 & 196 & \texttt{ACF00D823AAE\allowbreak{}8BC073F14ECF\allowbreak{}A322AED1B004\allowbreak{}CBBA85D5FA83\allowbreak{}416FEC4D24D8\allowbreak{}7482} & Frozen IEEE binary64 $J$, $J_g$, $b$, free map, and scales; transformed free coordinates; Euclidean norm. \\
M16 & 900 & \texttt{22A51CA5BB44\allowbreak{}03D1EB9499F4\allowbreak{}588D6EF19CE7\allowbreak{}D6C3636429D5\allowbreak{}EE2FCA6AE3AB\allowbreak{}207E} & Independently frozen IEEE binary64 $J$, $J_g$, $b$, free map, and scales; transformed free coordinates; Euclidean norm. \\
\bottomrule
\end{tabularx}


The M8 and M16 matrix snapshots were reproduced bitwise by independent capture
witnesses. Each application reconstructed every declared binary64 operand
exactly before high-precision comparison. The identity $A=J+E$ was checked
elementwise before either response route was evaluated.

\section{Finite-element route results}
\begin{table}[htbp]\centering\footnotesize
\caption{Rounded finite-element route values; exact decimal strings remain in the machine-readable decision files}\label{tab:S6}

\resizebox{\linewidth}{!}{%
\begin{tabular}{@{}llrrrrrl@{}}
\toprule
\textbf{Case} & \textbf{Dimension} & \textbf{Route} & \textbf{Observed response} & \textbf{Uncertainty bound} & \textbf{Lower bound} & \textbf{Rel.-error upper bound} & \textbf{Decision} \\
\midrule
M8 & 196 & Direct difference & $2.918501\times 10^{-12}$ & $5.880138\times 10^{-13}$ & $2.330488\times 10^{-12}$ & $2.523136\times 10^{-1}$ & Distinguishable \\
M8 & 196 & Response equation & $2.920014\times 10^{-12}$ & $2.588300\times 10^{-22}$ & $2.920014\times 10^{-12}$ & $8.863996\times 10^{-11}$ & Distinguishable \\
M16 & 900 & Direct difference & $2.463376\times 10^{-11}$ & $3.970800\times 10^{-12}$ & $2.066296\times 10^{-11}$ & $1.921699\times 10^{-1}$ & Distinguishable \\
M16 & 900 & Response equation & $2.463672\times 10^{-11}$ & $6.930483\times 10^{-21}$ & $2.463672\times 10^{-11}$ & $2.813071\times 10^{-10}$ & Distinguishable \\
\bottomrule
\end{tabular}%
}

\end{table}

\section{Finite-element precision, residual, identity, and repeatability}
\begin{table}[htbp]\centering\footnotesize
\caption{Rounded precision and identity checks; all gates passed}\label{tab:S7}
\begin{tabularx}{\linewidth}{>{\centering\arraybackslash}p{0.09\linewidth} *{4}{>{\centering\arraybackslash}X}}
\toprule
Case & 120/180 difference & Maximum residual & Route-identity difference & Repeat difference \\
\midrule
M8 & $8.81\times10^{-121}$ & $2.55\times10^{-124}$ & $9.18\times10^{-123}$ & $0$ \\
M16 & $3.97\times10^{-121}$ & $4.14\times10^{-124}$ & $3.04\times10^{-122}$ & $0$ \\
\bottomrule
\end{tabularx}
\end{table}

M8 required 468~s and 0.11~GiB recorded peak memory. M16 required 10.4~h and
1.35~GiB. These are offline reference costs for two fixed applications; they
are not a sparse-complexity fit or an online-estimator benchmark.

\section{Finite-element route-specific error components}
\begin{table}[htbp]\centering\footnotesize
\caption{Component-level coverage and final route enclosure}\label{tab:S8}
\begin{tabularx}{\linewidth}{>{\centering\arraybackslash}p{0.09\linewidth} >{\RaggedRight\arraybackslash}p{0.19\linewidth} X >{\centering\arraybackslash}p{0.16\linewidth} >{\centering\arraybackslash}p{0.12\linewidth}}
\toprule
Case & Route & Registered components & Final $U$ & Coverage \\
\midrule
M8 & Direct difference & two parent solves; subtraction & $5.880\times10^{-13}$ & 3/3 \\
M8 & Response equation & parent propagation; forcing; response solve & $2.588\times10^{-22}$ & 3/3 \\
M16 & Direct difference & two parent solves; subtraction & $3.971\times10^{-12}$ & 3/3 \\
M16 & Response equation & parent propagation; forcing; response solve & $6.930\times10^{-21}$ & 3/3 \\
\bottomrule
\end{tabularx}
\end{table}

All registered components covered their independently evaluated reference
errors for the two fixed applications. Coverage is operation-graph and case
specific; exact component strings remain in the archived ledgers.

\section{Historical diagnostics retained for provenance}
\begin{figure}[htbp]
\centering
\includegraphics[width=0.94\linewidth,height=0.68\textheight,keepaspectratio]{supplement_figures/F6_coverage_and_tightness.pdf}
\caption{Historical coverage and aggregate conservativeness retained for provenance. The main manuscript replaces this summary with factorized effectivity anatomy}\label{fig:S3}
\end{figure}

\begin{figure}[htbp]
\centering
\includegraphics[width=0.94\linewidth,height=0.68\textheight,keepaspectratio]{supplement_figures/S4_route_components.pdf}
\caption{Historical route-component anatomy for M8 and M16. Right-hand-side and solution-space components are not added across spaces}\label{fig:S4}
\end{figure}

\section{Claim boundaries and machine-readable ledger}
The CSV file \texttt{Paper2\_claim\_evidence\_matrix\_v0.10.0.csv} remains
the authoritative record for claims frozen through v0.10.0. The G2-20
execution report and adjudication extend it for the response-route theorem,
new holdout, independent backend, and effectivity analysis; historical claim
IDs and evidence tiers were not rewritten.
The complete row-level matrix is distributed as CSV rather than typeset at an
illegible scale in this PDF. Table~\ref{tab:S9} retains the concise inference
boundary needed for human review.

\begin{table}[htbp]\centering\small
\caption{Interpretation boundary by claim domain}\label{tab:S9}

\begin{tabularx}{\linewidth}{@{}>{\raggedright\arraybackslash}p{0.18\linewidth}>{\raggedright\arraybackslash}X>{\raggedright\arraybackslash}X@{}}
\toprule
\textbf{Domain} & \textbf{Established} & \textbf{Not established} \\
\midrule
Theory & Finite response identity, directional non-determination, and coordinate mapping under declared assumptions. & Prevalence, optimal scaling, or a universal guard rule. \\
Synthetic evidence & Pre-specified reproduction and coverage on the frozen development and disjoint holdout designs. & Unrestricted-system coverage or estimator tightness. \\
Frozen tangents & Route-specific qualification completed separately for M8 and M16. & Mesh law, dimension trend, or population-level generality. \\
Route comparison & Each route has its own uncertainty components and a case-specific decision. & Universal route ranking or superiority. \\
Interpretation & Numerical distinguishability under declared contracts. & Guard safety or harm, nonlinear consequence, finite-element stability, or engineering qualification. \\
Cost & Recorded offline cost and resource provenance for each frozen case. & Online estimator feasibility, complexity fit, or production scaling. \\
\bottomrule
\end{tabularx}

\end{table}

\section*{Data and code availability}
The existing non-confidential reproducibility artifact is archived as
\href{https://doi.org/10.5281/zenodo.21833850}{Zenodo concept DOI
10.5281/zenodo.21833850}; the public implementation is available from the
\href{https://github.com/yqma7980/scale-explicit-matrix-perturbation-observability}{GitHub repository}.
The v1.1.0 G2-20 release package contains the new holdout, route results,
independent-backend audit, analysis tables, scripts, and hashes. The archive
excludes Abaqus executables, user-element source, solver state histories, and
engineering project data.

\end{document}
